\section{関連研究}
\subsection{物体検出とオープンボキャブラリ}
YOLO\cite{ref_YOLO}とDETR\cite{ref_DETR}は，学習時に定めたクラス集合を前提とする物体検出モデルである．
CLIP\cite{ref_CLIP}は，画像とテキストの対の対照学習により，自然言語で指定した概念へのゼロショットの転移を可能にした．
OWL-ViT\cite{ref_OWLViT}とOWLv2\cite{ref_OWLv2}は，CLIPの画像エンコーダの各パッチからボックスを予測し，テキストで指定したカテゴリの物体を検出する．
OWLv2はクエリとは別に，クラス非依存の物体らしさのスコアも出力する．
Alpha-CLIP\cite{ref_AlphaCLIP}は，CLIPの画像エンコーダのパッチ埋め込みに注目領域を表すαチャネルの畳み込みを加え，画像全体の文脈を保ったまま指定した領域の特徴を求める．

\subsection{語彙と意味の階層}
VG\cite{ref_VG}は，10万枚を超える画像に物体を密に注釈し，物体の名前をWordNetのSynsetに正規化している．
前稿\cite{ref_prevVG}では，VGの注釈のSynsetの出現頻度から語彙を構成し，人手で定めた最上位の枠とWordNetの上下関係で木を作った．
さらに，全ての語の確率を部分木ごとに合計し，合計が閾値を超える最も深いノードを出力することで，COCO外の物体の命名を改善した．
下位クラスで分類してから親のクラスへ写像するCHiLS\cite{ref_CHiLS}も，具体的な語の確率を集約して粗い判断に用いる点で共通する．
階層のあるカテゴリの評価には，正解と出力の上位語の集合の重なりに基づく階層的適合率と階層的再現率\cite{ref_hier}を用いる．
本稿は，前稿の語彙の木による命名を，クラス非依存の検出器と領域を指定する分類器と組み合わせ，検出器と分類器の選択が命名の粒度と正しさに与える影響を調べる点が異なる．
